\section{“42”：哲学引用、命运与世界模型的边界}

结尾的“42”来自《银河系漫游指南》的梗：生命、宇宙以及一切问题的答案可能是 42。访谈最后，主持人问“这个世界是一个巨大世界模型吗”“你能预测命运吗”。谢赛宁回答，世界当然可以被看作一个巨大的世界模型，但我们不能预测命运，因为资源不够；也许要用地球或整个宇宙作为计算机，才能得到答案。

这段不是玩笑收尾，而是回到世界模型的边界：world model 是为了在有限资源下做足够好的预测，而不是全知全能地复制宇宙。模型越接近真实世界，越要面对计算资源、表征压缩、目标选择和不可知性的限制。

\subsection{放开维特根斯坦}

前面用“42”提醒世界模型不能假装全知；这里则回到语言和世界的关系。这个小节重要，是因为很多 LLM 叙事正是通过哲学名句把语言能力包装成世界理解。

谢赛宁对维特根斯坦引用的吐槽，实际上是在为世界模型划边界。早期维特根斯坦的语言-世界命题有特定哲学语境，后期又转向语言游戏。把“语言的边界就是世界的边界”直接拿来证明 LLM 能覆盖世界，是语境错配。

如果按后期“语言游戏”的思路，语言的意义来自实践和生活形式，那么这反而更接近世界模型观点：语言不是悬浮的 token，而是在真实世界行动、使用和反馈中获得意义。

\begin{importantbox}{语言与世界的关系}
语言可以压缩世界、沟通世界、组织经验，但不能替代世界。世界模型要学习的，是语言与实践发生关系的那部分结构：行动、反馈、约束、因果和可预测性。
\end{importantbox}

\subsection{The normal one 与团队电池}

最后回到 the normal one。谢赛宁借 Jürgen Klopp 的说法，把自己想象成团队里的电池：用热情和能量给别人发电。这不是鸡汤，而是 research organization 的一个核心功能。长期探索会有大量沮丧、失败和不确定，团队需要有人维持方向、信任和节奏。

他也承认研究底色常常悲凉：真正开心的时刻可能只有东西做出来的 5\% 到 10\%。这句话很真实，也解释了为什么世界模型这种方向需要组织文化支撑。没有长期心理韧性，宏大愿景会很快被短期挫败压垮。

\subsection{本章小结}

结尾把技术问题和人生问题压在一起：世界可以被建模，但模型永远受资源和目标限制；语言能启发思考，但不能替代实践；普通人也可以通过长期选择、团队能量和研究品味，进入大问题。

